\subsubsection{Trajectory Contribution Analysis}
\begin{enumerate}
    \item \textbf{Fortran Core Contribution Analysis Functions}
          
          The Fortran module provides core routines for trajectory contribution analysis
          using cosine similarity and spike detection.
          
          \begin{enumerate}
            \item \textbf{Trajectory Contribution Calculation}
                  
                  Calculates integrated trajectory contribution using cosine similarity between
                  factor and dependent vectors.
                  
                  \textbf{Mathematical Formula:}
                  \[
                    \text{contribution}= \frac{\text{factor} \cdot \text{dependent}}{\|\text{factor}\|
                        \cdot \|\text{dependent}\|}
                  \]
                  
                  \textbf{Arguments:}
                  \begin{itemize}
                    \item \texttt{n\_timepoints integer(int32)}: Number of timepoints
                          
                    \item \texttt{factor real(real64)}: Vector of independent variable [n\_timepoints]
                          
                    \item \texttt{dependent real(real64)}: Vector of dependent variable [n\_timepoints]
                          
                    \item \texttt{mode integer(int32)}: Processing mode (1 = Normal, 2 =
                          RAP)
                          
                    \item \texttt{contribution real(real64)}: Output contribution value
                          
                    \item \texttt{ierr integer(int32)}: Error code
                  \end{itemize}
                  \begin{lstlisting}[language=Fortran, caption={Trajectory Contribution Calculation}]
call trajectory_contribution(factor, dependent, n_timepoints, mode, &
                            contribution, ierr)
                  \end{lstlisting}
                  
            \item \textbf{Spike Contribution Calculation}
                  
                  Calculates timepoint-wise spike contributions using cosine similarity for
                  each individual timepoint.
                  
                  \textbf{Mathematical Formula:}
                  \[
                    \text{contribution}(i) = \frac{\text{factor}(i) \cdot \text{dependent}(i)}{\|\text{factor}\|
                        \cdot \|\text{dependent}\|}\quad \text{for }1 \le i \le n\_timepoints
                  \]
                  
                  \textbf{Arguments:}
                  \begin{itemize}
                    \item \texttt{n\_timepoints integer(int32)}: Number of timepoints
                          
                    \item \texttt{factor real(real64)}: Vector of independent variable [n\_timepoints]
                          
                    \item \texttt{dependent real(real64)}: Vector of dependent variable [n\_timepoints]
                          
                    \item \texttt{mode integer(int32)}: Processing mode (1 = Normal, 2 =
                          RAP)
                          
                    \item \texttt{contribution real(real64)}: Output spike contributions
                          [n\_timepoints]
                          
                    \item \texttt{ierr integer(int32)}: Error code
                  \end{itemize}
                  \begin{lstlisting}[language=Fortran, caption={Spike Contribution Calculation}]
call spike_contribution(factor, dependent, n_timepoints, mode, &
                        contribution, ierr)
                  \end{lstlisting}
                  
            \item \textbf{Calculate Contributions (Expert Version)}
                  
                  Calculates both spike and integrated contributions for a specific factor
                  with pre-allocated work arrays.
                  
                  \textbf{Arguments:}
                  \begin{itemize}
                    \item \texttt{n\_factors integer(int32)}: Number of factors in trajectories
                          
                    \item \texttt{n\_samples integer(int32)}: Number of samples in trajectories
                          
                    \item \texttt{n\_timepoints integer(int32)}: Number of timepoints in
                          trajectories
                          
                    \item \texttt{trajectories real(real64)}: 3D trajectories array [n\_factors,
                    n\_samples, n\_timepoints]
                    
                    \item \texttt{i\_factor integer(int32)}: Factor index to calculate contributions
                          for
                          
                    \item \texttt{dependent\_idx integer(int32)}: Dependent variable index
                          
                    \item \texttt{mode integer(int32)}: Processing mode (1 = Normal, 2 =
                          RAP)
                          
                    \item \texttt{integrated\_contribs real(real64)}: Output integrated contributions
                          [n\_samples]
                          
                    \item \texttt{spike\_contribs real(real64)}: Output spike contributions
                          [n\_timepoints, n\_samples]
                          
                    \item \texttt{temp\_factor\_vector real(real64)}: Work array [n\_timepoints]
                          
                    \item \texttt{temp\_dependent\_vector real(real64)}: Work array [n\_timepoints]
                          
                    \item \texttt{ierr integer(int32)}: Error code
                  \end{itemize}
                  \begin{lstlisting}[language=Fortran, caption={Expert Contributions Calculation}]
call calc_contributions(trajectories, n_factors, n_samples, & 
                n_timepoints, i_factor, dependent_idx, mode, &
                spike_contribs, integrated_contribs, &
                temp_factor_vector, temp_dependent_vector, ierr)
                  \end{lstlisting}
                  
            \item \textbf{Calculate Contributions (Allocation Version)}
                  
                  Calculates both spike and integrated contributions with internal memory
                  allocation.
                  
                  \textbf{Arguments:}
                  \begin{itemize}
                    \item \texttt{n\_factors integer(int32)}: Number of factors in trajectories
                          
                    \item \texttt{n\_samples integer(int32)}: Number of samples in trajectories
                          
                    \item \texttt{n\_timepoints integer(int32)}: Number of timepoints in
                          trajectories
                          
                    \item \texttt{trajectories real(real64)}: 3D trajectories array [n\_factors,
                    n\_samples, n\_timepoints]
                    
                    \item \texttt{i\_factor integer(int32)}: Factor index to calculate contributions
                          for
                          
                    \item \texttt{dependent\_idx integer(int32)}: Dependent variable index
                          
                    \item \texttt{mode integer(int32)}: Processing mode (1 = Normal, 2 =
                          RAP)
                          
                    \item \texttt{integrated\_contribs real(real64)}: Output integrated contributions
                          [n\_samples]
                          
                    \item \texttt{spike\_contribs real(real64)}: Output spike contributions
                          [n\_timepoints, n\_samples]
                          
                    \item \texttt{ierr integer(int32)}: Error code
                  \end{itemize}
                  \begin{lstlisting}[language=Fortran, caption={Allocation-Based Contributions Calculation}]
call calc_contributions_alloc(trajectories, n_factors, n_samples, &
                    n_timepoints, i_factor, dependent_idx, mode, &
                    spike_contribs, integrated_contribs, ierr)
                  \end{lstlisting}
                  
            \item \textbf{Vector Accessor Routines}
                  
                  Extract specific slices from the 3D trajectories array.
                  
                  \textbf{Arguments for Sample Extraction:}
                  \begin{itemize}
                    \item \texttt{n\_factors integer(int32)}: Number of factors
                          
                    \item \texttt{n\_samples integer(int32)}: Number of samples
                          
                    \item \texttt{n\_timepoints integer(int32)}: Number of timepoints
                          
                    \item \texttt{trajectories real(real64)}: 3D trajectories array
                          
                    \item \texttt{i\_factor integer(int32)}: Factor index
                          
                    \item \texttt{i\_timepoint integer(int32)}: Timepoint index
                          
                    \item \texttt{result\_vector real(real64)}: Output vector [n\_samples]
                          
                    \item \texttt{ierr integer(int32)}: Error code
                  \end{itemize}
                  \begin{lstlisting}[language=Fortran, caption={Extract Samples Vector}]
call get_vec_across_samples(trajectories, n_factors, n_samples, &
                        n_timepoints, i_factor, i_timepoint, &
                        result_vector, ierr)
                  \end{lstlisting}
                  
                  \textbf{Arguments for Timepoint Extraction:}
                  \begin{itemize}
                    \item \texttt{n\_factors integer(int32)}: Number of factors
                          
                    \item \texttt{n\_samples integer(int32)}: Number of samples
                          
                    \item \texttt{n\_timepoints integer(int32)}: Number of timepoints
                          
                    \item \texttt{trajectories real(real64)}: 3D trajectories array
                          
                    \item \texttt{i\_factor integer(int32)}: Factor index
                          
                    \item \texttt{i\_sample integer(int32)}: Sample index
                          
                    \item \texttt{result\_vector real(real64)}: Output vector [n\_timepoints]
                          
                    \item \texttt{ierr integer(int32)}: Error code
                  \end{itemize}
                  \begin{lstlisting}[language=Fortran, caption={Extract Timepoints Vector}]
call get_vec_across_timepoints(trajectories, n_factors, n_samples, &
                            n_timepoints, i_factor, i_sample, &
                            result_vector, ierr)
                  \end{lstlisting}
                  
            \item \textbf{Outlier detection}
                  
                  The outlier detection workflow consists of two main pathways:
                  integrated (trajectory-level) analysis and spike (timepoint-level)
                  analysis. You can perform outlier detection by calling the following subroutines
                  in sequence:
                  
                  \begin{enumerate}
                    \item \textbf{Calculate Integrated Threshold}
                          
                          The subroutine calculates an empirical threshold for integrated (trajectory-level)
                          contributions using pre-computed permutations. This determines
                          which samples are significant overall across all genes. Call the \texttt{calc\_integrated\_threshold}
                          subroutine with the following arguments:
                          
                          \textbf{Input arguments:}
                          \begin{itemize}
                            \item \texttt{n\_samples integer(int32)}: Number of samples
                                  
                            \item \texttt{contributions real(real64)}: 1D array of integrated
                                  contributions (size: \texttt{n\_samples})
                                  
                            \item \texttt{percentile\_val real(real64)}: Percentile value for
                                  threshold (0.0--100.0)
                                  
                            \item \texttt{permutation integer(int32)}: Pre-computed permutation
                                  indices for sorted contributions (size: \texttt{n\_samples})
                          \end{itemize}
                          \textbf{Output arguments:}
                          \begin{itemize}
                            \item \texttt{threshold real(real64)}: Scalar threshold value
                                  
                            \item \texttt{ierr integer(int32)}: Error code (for error
                                  details please check
                                  \hyperref[sec:error-handling]{error handling section})
                          \end{itemize}
                          \begin{lstlisting}[language=Fortran, caption={Integrated Threshold Calculation}]
call calc_integrated_threshold(contributions, n_samples, &
                            percentile_val, threshold, &
                            permutation, ierr)
                          \end{lstlisting}
                          
                    \item \textbf{Detect Integrated Outliers}
                          
                          The subroutine identifies outlier samples where the integrated contribution
                          exceeds the empirical threshold. Used to find trajectories with
                          significant overall explanatory contributions. Call the \texttt{detect\_outliers\_integrated}
                          subroutine with the following arguments:
                          
                          \textbf{Input arguments:}
                          \begin{itemize}
                            \item \texttt{n\_samples integer(int32)}: Number of samples
                                  
                            \item \texttt{contributions real(real64)}: Array of integrated contributions
                                  (size: \texttt{n\_samples})
                                  
                            \item \texttt{threshold real(real64)}: Scalar threshold value for
                                  outlier detection
                          \end{itemize}
                          \textbf{Output arguments:}
                          \begin{itemize}
                            \item \texttt{outlier\_mask logical}: 1D logical array indicating
                                  outliers (size: \texttt{n\_samples})
                                  
                            \item \texttt{ierr integer(int32)}: Error code (for error
                                  details please check
                                  \hyperref[sec:error-handling]{error handling section})
                          \end{itemize}
                          \begin{lstlisting}[language=Fortran, caption={Integrated Outlier Detection}]
call detect_outliers_integrated(contributions, n_samples,&
                                threshold, outlier_mask, ierr)
                          \end{lstlisting}
                          
                    \item \textbf{Calculate Spike Thresholds}
                          
                          The subroutine calculates empirical thresholds for spike contributions
                          using pre-sorted permutations. Computes percentile-based
                          thresholds for each timepoint in spike contribution data. Call the
                          \texttt{calc\_spike\_thresholds} subroutine with the following
                          arguments:
                          
                          \textbf{Input arguments:}
                          \begin{itemize}
                            \item \texttt{n\_timepoints integer(int32)}: Number of timepoints
                                  
                            \item \texttt{n\_samples integer(int32)}: Number of samples
                                  
                            \item \texttt{spike\_contribs real(real64)}: 2D array of spike contributions
                                  (size: \texttt{n\_timepoints} × \texttt{n\_samples})
                                  
                            \item \texttt{percentile\_val real(real64)}: Percentile value for
                                  threshold (0.0--100.0)
                                  
                            \item \texttt{permutation integer(int32)}: Pre-computed permutation
                                  indices (size: \texttt{n\_samples} × \texttt{n\_timepoints}).
                          \end{itemize}
                          \textbf{Output arguments:}
                          \begin{itemize}
                            \item \texttt{thresholds real(real64)}: 1D array of thresholds for
                                  each timepoint (size: \texttt{n\_timepoints})
                                  
                            \item \texttt{ierr integer(int32)}: Error code (for error
                                  details please check
                                  \hyperref[sec:error-handling]{error handling section})
                          \end{itemize}
                          \begin{lstlisting}[language=Fortran, caption={Spike Thresholds Calculation}]
call calc_spike_thresholds(spike_contribs, n_timepoints, &
                        n_samples, percentile_val, thresholds, & 
                        permutation, ierr)
                          \end{lstlisting}
                          
                    \item \textbf{Detect Spike Outliers}
                          
                          The subroutine identifies outliers in spike contribution data by comparing
                          against timepoint-specific thresholds. Performs element-wise
                          comparison across all timepoints and samples. Call the \texttt{detect\_outliers\_spike}
                          subroutine with the following arguments:
                          
                          \textbf{Input arguments:}
                          \begin{itemize}
                            \item \texttt{n\_timepoints integer(int32)}: Number of timepoints
                                  
                            \item \texttt{n\_samples integer(int32)}: Number of samples
                                  
                            \item \texttt{spike\_contribs real(real64)}: 2D array of spike contributions
                                  (size: \texttt{n\_timepoints} × \texttt{n\_samples})
                                  
                            \item \texttt{thresholds real(real64)}: Array of thresholds for each
                                  timepoint (size: \texttt{n\_timepoints})
                          \end{itemize}
                          \textbf{Output arguments:}
                          \begin{itemize}
                            \item \texttt{outlier\_mask logical}: 2D logical array indicating
                                  outliers (size: \texttt{n\_timepoints} × \texttt{n\_samples})
                                  
                            \item \texttt{ierr integer(int32)}: Error code (for error
                                  details please check
                                  \hyperref[sec:error-handling]{error handling section})
                          \end{itemize}
                          \begin{lstlisting}[language=Fortran, caption={Spike Outlier Detection}]
call detect_outliers_spike(spike_contribs, n_timepoints, &
                        n_samples, thresholds, outlier_mask, ierr)
                          \end{lstlisting}
                  \end{enumerate}
                  
                  \textbf{For convenience, allocation-friendly versions are also available:}
                  \begin{itemize}
                    \item \texttt{calc\_integrated\_threshold\_alloc}: Handles internal allocation
                          and sorting
                          
                    \item \texttt{calc\_spike\_thresholds\_alloc}: Handles internal allocation
                          and sorting
                  \end{itemize}
                  
            \item \textbf{Process Trajectories with Per-Timepoint Percentiles}
                  
                  Complete pipeline processing with timepoint-specific spike thresholds.
                  
                  \textbf{Arguments:}
                  \begin{itemize}
                    \item \texttt{n\_factors integer(int32)}: Number of factors
                          
                    \item \texttt{n\_samples integer(int32)}: Number of samples
                          
                    \item \texttt{n\_timepoints integer(int32)}: Number of timepoints
                          
                    \item \texttt{factor\_mask\_count integer(int32)}: Number of true values
                          in factor mask
                          
                    \item \texttt{trajectories real(real64)}: 3D trajectories array
                          
                    \item \texttt{factor\_mask logical}: Mask for included factors [n\_factors]
                          
                    \item \texttt{dependent\_idx integer(int32)}: Dependent variable index
                          
                    \item \texttt{mode integer(int32)}: Processing mode (1 = Normal, 2 =
                          RAP)
                          
                    \item \texttt{percentile real(real64)}: Percentile value (0.0--100.0)
                          
                    \item \texttt{integrated\_contribs real(real64)}: Output integrated contributions
                          [n\_samples, factor\_mask\_count]
                          
                    \item \texttt{spike\_contribs real(real64)}: Output spike contributions
                          [n\_timepoints, n\_samples, factor\_mask\_count]
                          
                    \item \texttt{thresholds\_integrated\_contrib real(real64)}: Integrated
                          thresholds [factor\_mask\_count]
                          
                    \item \texttt{thresholds\_spike\_contrib real(real64)}: Spike thresholds
                          [n\_timepoints, factor\_mask\_count]
                          
                    \item \texttt{outliers\_integrated\_contrib logical}: Integrated outliers
                          [n\_samples, factor\_mask\_count]
                          
                    \item \texttt{outliers\_spike\_contrib logical}: Spike outliers [n\_timepoints,
                    n\_samples, factor\_mask\_count]
                    
                    \item \texttt{ierr integer(int32)}: Error code
                  \end{itemize}
                  \begin{lstlisting}[language=Fortran, caption={Complete Pipeline with Per-Timepoint Thresholds}]
call process_trajectories_alloc(trajectories, n_factors, n_samples, &
                    n_timepoints, factor_mask, factor_mask_count, & 
                    dependent_idx, mode, percentile, &
                    integrated_contribs,spike_contribs, &
                    thresholds_integrated_contrib, &
                    outliers_integrated_contrib, &
                    thresholds_spike_contrib, &
                    outliers_spike_contrib, ierr)
                  \end{lstlisting}
                  
            \item \textbf{Process Trajectories with Global Percentile}
                  
                  Complete pipeline processing with global spike thresholds (flattened
                  approach).
                  
                  \textbf{Arguments:}
                  \begin{itemize}
                    \item \texttt{n\_factors integer(int32)}: Number of factors
                          
                    \item \texttt{n\_samples integer(int32)}: Number of samples
                          
                    \item \texttt{n\_timepoints integer(int32)}: Number of timepoints
                          
                    \item \texttt{factor\_mask\_count integer(int32)}: Number of true values
                          in factor mask
                          
                    \item \texttt{trajectories real(real64)}: 3D trajectories array
                          
                    \item \texttt{factor\_mask logical}: Mask for included factors [n\_factors]
                          
                    \item \texttt{dependent\_idx integer(int32)}: Dependent variable index
                          
                    \item \texttt{mode integer(int32)}: Processing mode (1 = Normal, 2 =
                          RAP)
                          
                    \item \texttt{percentile real(real64)}: Percentile value (0.0--100.0)
                          
                    \item \texttt{integrated\_contribs real(real64)}: Output integrated contributions
                          [n\_samples, factor\_mask\_count]
                          
                    \item \texttt{spike\_contribs real(real64)}: Output spike contributions
                          [n\_timepoints, n\_samples, factor\_mask\_count]
                          
                    \item \texttt{thresholds\_integrated\_contrib real(real64)}: Integrated
                          thresholds [factor\_mask\_count]
                          
                    \item \texttt{thresholds\_spike\_contrib real(real64)}: Spike thresholds
                          [factor\_mask\_count] (1D)
                          
                    \item \texttt{outliers\_integrated\_contrib logical}: Integrated outliers
                          [n\_samples, factor\_mask\_count]
                          
                    \item \texttt{outliers\_spike\_contrib logical}: Spike outliers [n\_timepoints,
                    n\_samples, factor\_mask\_count]
                    
                    \item \texttt{ierr integer(int32)}: Error code
                  \end{itemize}
                  \begin{lstlisting}[language=Fortran, caption={Complete Pipeline with Global Threshold}]
call process_trajectories_flat_alloc(trajectories, n_factors, & 
                    n_samples, n_timepoints, factor_mask, &
                    factor_mask_count, dependent_idx, mode, &
                    percentile, integrated_contribs, &
                    spike_contribs, thresholds_integrated_contrib, &
                    outliers_integrated_contrib, &
                    thresholds_spike_contrib, &
                    outliers_spike_contrib, ierr)
                  \end{lstlisting}
          \end{enumerate}
          
          \textbf{Processing Modes:}
          \begin{itemize}
            \item \texttt{MODE\_NORMAL = 1}: Returns cosine similarity for normal vectors
                  
            \item \texttt{MODE\_RAP = 2}: Returns arccos of cosine similarity for RAP
                  projected vectors
          \end{itemize}
          
          \textbf{Key Features:}
          \begin{itemize}
            \item \textbf{Pure Subroutines}: Most routines are pure with no side effects
                  
            \item \textbf{Robust Error Handling}: Comprehensive error checking and validation
                  
            \item \textbf{Memory Efficient}: Expert versions allow pre-allocated work
                  arrays
                  
            \item \textbf{Mathematically Sound}: Uses cosine similarity with proper normalization
                  
            \item \textbf{Flexible Processing}: Supports both per-timepoint and global
                  thresholding
          \end{itemize}
\end{enumerate}